
\subsubsection{4.1 Dataset}

The Unterthiner MNIST zoo was used: 29,997 trained 4-layer convolutional neural networks with per-neuron weight vectors reshaped to token format (fan\_in = 16 per layer, 4 layers). The original Unterthiner FC-MLP zoo URL was unavailable (HTTP 404); the CNN zoo was adapted to maintain the permutation structure. Data was enriched with train\_acc from metrics.csv.gz to enable generalization gap computation (train\_loss - test\_loss).

\begin{table}[htbp]
\centering
\resizebox{\columnwidth}{!}{%
\begin{tabular}{ll}
\toprule
Property & Value \\
\midrule
Total models & 29,997 \\
Training split & 23,997 (80\%) \\
Test split & 6,000 (20\%) \\
Network depth & 4 layers \\
Fan-in per layer & 16 \\
Prediction target & Generalization gap \\
Split seed & 42 \\
\bottomrule
\end{tabular}
}
\end{table}

\subsubsection{4.2 Training Protocol}

All encoders were trained with Adam optimizer (lr = 0.001, beta = (0.9, 0.999), weight decay = 0.0001), CosineAnnealingLR scheduler (T\_max = 100, eta\_min = 1e-5), batch size 64. For the primary comparison (h-e1), training ran for 50 epochs with seed 42. For the six-encoder ablation (h-m1), training ran for 100 epochs with seeds {42, 123, 456}, yielding 18 training runs (6 encoders x 3 seeds). The h-m2 evaluation reused h-m1 checkpoints without retraining.

\subsubsection{4.3 Hardware}

Experiments ran on a single NVIDIA H100 NVL GPU. Total: 21 training runs (h-e1: 2 runs, h-m1: 18 runs, h-m2: evaluation-only).
